%! Exponential Function Manipulation — Unit 4, Lesson 4.4 — Group Activity (Answer Key)
\documentclass[10pt]{article}
\usepackage{precalculus-article}
\usepackage{precalculus-key}   % pulls in -boxes; do NOT also load -boxes
\newcommand{\TallMath}[1]{$\displaystyle #1\rule[-1.4em]{0pt}{3.2em}$}

\begin{document}

\pageheader{Unit 4, Lesson 4.4}{Group Activity --- Answer Key}
\namepartnerperiod

\vspace{0.1in}

\begin{scenariobox}[Equivalent Forms of Exponential Functions]{plumlight}
\small
Every exponential function can be written in many equivalent forms. Choosing the right form
makes certain features of the function immediately visible. In this activity you will practice
rewriting exponential expressions using the product property, the power property, and the
rules for negative and fractional exponents.
\end{scenariobox}

\vspace{0.1in}

\begin{tcolorbox}[colback=white, colframe=black!40,
    title=\textbf{Tier R --- Remediate},
    fonttitle=\bfseries, arc=2mm, left=3mm, right=3mm, top=2mm, bottom=2mm]

For each expression, (i) identify which property to apply and (ii) rewrite in the form $ab^x$.
Show each step.

\renewcommand{\arraystretch}{2.4}
\begin{tabularx}{\linewidth}{@{} >{\centering\arraybackslash}p{2.4cm} >{\centering\arraybackslash}p{3.2cm} X X @{}}
\textbf{Expression} & \textbf{Property} & \textbf{Step 1} & \textbf{Final: $ab^x$} \\
\midrule
$2^{x+4}$  & \ans{product} & $2^x \cdot 2^{{\color{keyred}\mathbf{4}}}$ & \ans{$16 \cdot 2^x$} \\
$3^{2x}$   & \ans{power}   & $(3^{{\color{keyred}\mathbf{2}}})^x$       & \ans{$9^x$} \\
$5^{-x}$   & \ans{negative exp.} & $\left(\dfrac{1}{{\color{keyred}\mathbf{5}}}\right)^x$ & \ans{$\left(\tfrac{1}{5}\right)^x$} \\
$4^{x/2}$  & \ans{power}   & $(4^{{\color{keyred}\mathbf{1/2}}})^x$     & \ans{$2^x$} \\
\end{tabularx}
\end{tcolorbox}

\vspace{0.1in}

\begin{tcolorbox}[colback=white, colframe=black!40,
    title=\textbf{Tier A --- Apply},
    fonttitle=\bfseries, arc=2mm, left=3mm, right=3mm, top=2mm, bottom=2mm]

\begin{enumerate}[itemsep=12pt]

  \item Three functions are given below. Show that all three are equivalent by rewriting each
        in the form $a \cdot b^x$.

        \[
          f(x) = 4^{x+1}
          \qquad
          g(x) = 4 \cdot 4^x
          \qquad
          h(x) = 16 \cdot \left(\tfrac{1}{4}\right)^{-x+1}
        \]

        $f(x) = $ \ans{$4 \cdot 4^x$} \quad $g(x) = $ \ans{$4 \cdot 4^x$} (already in form)

        $h(x)$: $\left(\tfrac{1}{4}\right)^{-x} = 4^x$, so $h(x) = 16 \cdot \left(\tfrac{1}{4}\right)^{-x+1} = 16 \cdot \tfrac{1}{4} \cdot 4^x = $ \ans{$4 \cdot 4^x$}

        All three simplify to: \ans{$4 \cdot 4^x$}

  \item Rewrite $f(x) = 8^{x/3}$ as a single base raised to the $x$ power. Show your work.

        \vspace{0.04in}
        $8^{x/3} = (8^{{\color{keyred}\mathbf{1/3}}})^x = $ \ans{$2^x$}

  \item Rewrite $g(x) = 27 \cdot 3^{-x}$ in the form $a \cdot b^x$. Classify as growth or decay
        and state the $y$-intercept.

        \vspace{0.04in}
        $g(x) = 27 \cdot 3^{-x} = 27 \cdot \left(\dfrac{1}{3}\right)^x = $ \ans{$27\cdot\left(\dfrac{1}{3}\right)^x$}

        \vspace{0.04in}
        Growth / Decay (circle one): \ans{Decay} \quad $y$-intercept: \ans{$(0,27)$}

\end{enumerate}
\end{tcolorbox}

\vspace{0.1in}

\begin{tcolorbox}[colback=white, colframe=black!40,
    title=\textbf{Tier E --- Extend},
    fonttitle=\bfseries, arc=2mm, left=3mm, right=3mm, top=2mm, bottom=2mm]

\begin{enumerate}[itemsep=10pt]

  \item Starting with $f(x) = 6^{2x+1}$, produce \textbf{three} different equivalent forms.

        \textbf{Form 1} (vertical dilation): $6^{2x+1} = 6 \cdot 6^{2x}$ \ans{$= 6 \cdot 36^x$}

        \textbf{Form 2} (single base to the $x$ power): $6^{2x} = (6^2)^x$, so $f(x) = 6 \cdot (6^2)^x$ \ans{$= 6 \cdot 36^x$}

        \textbf{Form 3} (creative): \ans{$f(x) = 6^{2x+1} = (36)^x \cdot 6$, or write as $6 \cdot 36^x$ emphasizing $a = 6, b = 36$.
        Another option: $(1296)^{x/2} \cdot 6$.}

        \begin{teachernote}
          Accept any valid equivalent form for Form 3. Check that students use $f(x) = 6 \cdot 36^x$
          (most natural $ab^x$ form) and clearly justify their algebra.
        \end{teachernote}

  \item For each context, which form is most useful?

        \textbf{Context A} ($y$-intercept): \ans{Form 1 ($6 \cdot 36^x$), because $a = 6$ is immediately visible as $f(0)$.}

        \textbf{Context B} (compare to $36^x$): \ans{Form 2 or Form 1 --- writing the base as $36$ makes the comparison direct.}

        \textbf{Context C} (horizontal stretch from $6^x$): \ans{The original form $6^{2x+1}$ shows the coefficient $c = 2$ inside the exponent most clearly.}

\end{enumerate}
\end{tcolorbox}

\end{document}
